\chapter{The CMG Correspondence}\label{ch:cmg}

\begin{assumption}\label{ass:ch13}
Throughout this chapter $\cC$ is a locally small category with an initial
object $0$ and a terminal object $1$. Small coproducts in $\cC$ are not assumed
unless stated. We write $\cC^\conn$ for the full subcategory of connected
objects (\cref{def:connected}).
\end{assumption}

This chapter proves the CMG Correspondence Theorem, which relates three notions
introduced in earlier chapters:
\begin{itemize}
\item the \emph{Chloe Condition} on a cat (\cref{def:chloe}): every leg is
  connected;
\item \emph{Margaret Compatibility} of a cocone $\beta\colon K\to\iota X$
  (\cref{def:margaret}): every morphism from a connected object to $X$ factors
  uniquely through exactly one leg;
\item \emph{Gideon Decompositions} of an object $X$ (\cref{def:gideon}): a cat
  satisfying the Chloe Condition together with a coproduct cocone from its legs
  to $X$.
\end{itemize}
We then deduce a recognition theorem for categories of cats and study the
interaction of the three notions with the monoidal structure of
\cref{ch:monoidal}. The identification of each notion with a familiar one is
recorded in \cref{rem:cmg-honest}.

\section{The CMG Correspondence Theorem}\label{sec:cmg-theorem}

\begin{theorem}[CMG Correspondence]\label{thm:cmg}
Let $X$ be an object of $\cC$, let $K$ be a $\cC$-cat, and let $\beta\colon K\to\iota X$
be a cocone. Consider the conditions:
\begin{itemize}
\item[\textup{(C)}] $K$ satisfies the Chloe Condition;
\item[\textup{(M)}] $\beta$ is Margaret compatible;
\item[\textup{(G)}] $(K,\beta)$ is a Gideon Decomposition of $X$.
\end{itemize}
Then:
\begin{enumerate}
\item\label{it:cmg-forward} \textup{(G)} implies \textup{(C)} and \textup{(M)}.
\item\label{it:cmg-backward} Suppose that $X$ admits a Gideon Decomposition
  $(K_0,\beta_0)$. If \textup{(C)} and \textup{(M)} hold, then there is a unique
  morphism of cats $\theta\colon K\to K_0$ with $\beta_0\theta=\beta$; this
  $\theta$ is an isomorphism, and \textup{(G)} holds.
\item\label{it:cmg-unique} The category $\mathsf{Gid}(X)$ is either empty or
  equivalent to $\one$: any two Gideon Decompositions of $X$ are related by a
  unique morphism, and it is an isomorphism.
\end{enumerate}
\end{theorem}

\begin{proof}
\ref{it:cmg-forward} (C) is part of the definition of a Gideon Decomposition.
For (M), let $Y$ be connected. Since $X$, with the coprojections $\beta_\alpha$,
is a coproduct of the family $(L_\alpha)$, the map $m^\beta_Y$ is precisely the
canonical map \eqref{eq:connected} for this family, which is bijective because
$Y$ is connected.

\ref{it:cmg-backward} By \ref{it:cmg-forward}, $K_0$ satisfies (C) and $\beta_0$
satisfies (M). By \cref{prop:margaret-equiv}, $\Yon^\conn(\beta)\colon\Yon^\conn K\to y^\conn X$
and $\Yon^\conn(\beta_0)\colon\Yon^\conn K_0\to y^\conn X$ are isomorphisms. Since $K$
satisfies the Chloe Condition, \cref{prop:restricted-yoneda} shows that
$\Yon^\conn\colon\Fel(\cC)(K,K_0)\to\PSh(\cC^\conn)(\Yon^\conn K,\Yon^\conn K_0)$ is
a bijection; let $\theta$ be the unique morphism with
$\Yon^\conn(\theta)=\Yon^\conn(\beta_0)^{-1}\circ\Yon^\conn(\beta)$. Applying the same
proposition to the pair $(K,\iota X)$, the map
$\Yon^\conn\colon\Fel(\cC)(K,\iota X)\to\PSh(\cC^\conn)(\Yon^\conn K,y^\conn X)$ is a
bijection. Since $\Yon^\conn(\beta_0\theta)=\Yon^\conn(\beta)$, we obtain
$\beta_0\theta=\beta$. If $\theta'$ also satisfies $\beta_0\theta'=\beta$, then
$\Yon^\conn(\beta_0)\Yon^\conn(\theta')=\Yon^\conn(\beta)$, so
$\Yon^\conn(\theta')=\Yon^\conn(\theta)$ and $\theta'=\theta$. Since $K$ and
$K_0$ both satisfy the Chloe Condition, \cref{prop:restricted-yoneda} shows that
$\Yon^\conn$ is fully faithful on the full subcategory spanned by $K$ and $K_0$,
and so reflects isomorphisms there; as $\Yon^\conn(\theta)$ is an isomorphism,
so is $\theta$. Finally, write $\theta=(\varphi,(\theta_\alpha))$. By
\cref{prop:iso}, $\varphi$ is bijective and each $\theta_\alpha$ is invertible, so
the family $\beta_\alpha=\beta_{0,\varphi(\alpha)}\circ\theta_\alpha$ is a coproduct cocone,
being obtained from the coproduct cocone $\beta_0$ by reindexing along a
bijection and precomposing with isomorphisms. Hence (G) holds.

\ref{it:cmg-unique} Given Gideon Decompositions $(K,\beta)$ and $(K_0,\beta_0)$,
part \ref{it:cmg-forward} shows that $(K,\beta)$ satisfies (C) and (M), and part
\ref{it:cmg-backward} provides a unique morphism $(K,\beta)\to(K_0,\beta_0)$,
which is an isomorphism.
\end{proof}

\begin{corollary}\label{cor:cmg-iff}
If $X$ is Gideon-decomposable, then a cocone $\beta\colon K\to\iota X$ from a cat
satisfying the Chloe Condition is a Gideon Decomposition if and only if it is
Margaret compatible. In symbols: under the hypothesis that $\mathsf{Gid}(X)\neq\varnothing$,
\[
\textup{(G)}\iff\textup{(C)}\wedge\textup{(M)}.
\]
\end{corollary}

\begin{proof}
Combine \cref{thm:cmg}\ref{it:cmg-forward},\ref{it:cmg-backward}.
\end{proof}

\subsection*{Sharpness}

Each hypothesis in \cref{thm:cmg,cor:cmg-iff} is necessary.

\begin{counterexample}[Existence cannot be dropped]\label{cex:cmg-existence}
In each of the following, the object $X$ admits no Gideon Decomposition.
\begin{enumerate}
\item In $\Top$, the cocone of \cref{ex:margaret}(iii) from the cat
  $K=\langle\{q\}\rangle_{q\in\mathbb Q}$ to $\mathbb Q$ satisfies (C) and (M) but
  not (G). Here $\mathbb Q$ admits no Gideon Decomposition (\cref{ex:gideon}(ii)).
\item In $\Ab$, for any nonzero abelian group $X$, the unique cocone from the
  legless cat to $X$ satisfies (C) vacuously and (M) vacuously
  (\cref{ex:margaret}(ii)), but not (G), since $X$ is not initial.
\end{enumerate}
\end{counterexample}

\begin{counterexample}[The Chloe Condition cannot be dropped]\label{cex:cmg-chloe}
In $\Set$ let $Y$ be a set with at least one element and $X=Y\sqcup Y$. The
unipedal cat $\iota X$ with the identity cocone, and the two-legged cat
$\langle Y,Y\rangle$ with the coprojections, are both Margaret compatible cocones
exhibiting $X$ as a coproduct of their legs. They are not isomorphic over $X$,
since their leg sets have different cardinalities. Without (C), therefore,
neither the implication \textup{(M)}$\Rightarrow$\textup{(G)} nor the uniqueness
of \cref{thm:cmg}\ref{it:cmg-unique} survives. In this example both cocones
are coproduct cocones; the following shows that without (C), Margaret
Compatibility does not even force a coproduct cocone, although $X$ is
Gideon-decomposable. In $\Ab$, the object $X=0$ has the legless Gideon
Decomposition. The cocone from $\iota(\mathbb Z)$ to $0$ is Margaret compatible,
vacuously, since $\Ab$ has no connected objects, but it is not a coproduct
cocone, since $\mathbb Z$ is not a zero object.
\end{counterexample}

\begin{counterexample}[Margaret Compatibility cannot be dropped]\label{cex:cmg-margaret}
In $\Set$, the unipedal cat $\iota(1)$ with the cocone $1\to2$ picking out one
element satisfies (C), and $2$ is Gideon-decomposable, but the cocone is
neither Margaret compatible nor a Gideon Decomposition.
\end{counterexample}

\begin{remark}\label{rem:cmg-honest}
We identify the three notions and the theorem with classical ones.
\begin{enumerate}
\item The Chloe Condition on a cat is the condition that its legs be connected
  objects; Chloe cats are, up to equivalence, families of connected objects
  (\cref{prop:chloe-basic}(iv)).
\item Margaret Compatibility of a cocone is the invertibility of the induced
  map of restricted Yoneda presheaves on connected objects
  (\cref{prop:margaret-equiv}).
\item A Gideon Decomposition is a decomposition of an object as a coproduct of
  connected objects, i.e.\ a decomposition into connected components.
\item \Cref{thm:cmg}\ref{it:cmg-unique} is the uniqueness of decompositions into
  connected components, and \cref{thm:cmg}\ref{it:cmg-backward} is the
  observation that, between coproducts of connected objects, a morphism is
  determined by, and is invertible if and only if it induces a bijection on,
  morphisms out of connected objects.
\end{enumerate}
We are not aware of any application of the CMG Correspondence Theorem outside
this monograph. Within it, the theorem is used only to show that the Gideon
number (\cref{def:gideon-number}) is well defined and in the proof of
\cref{thm:recognition}.
\end{remark}

\begin{definition}\label{def:gideon-number}
For a Gideon-decomposable object $X$, its \emph{Gideon number} $\gamma(X)$ is the
cardinality of the leg set of any Gideon Decomposition of $X$. It is well
defined by \cref{thm:cmg}\ref{it:cmg-unique} and \cref{prop:iso}.
\end{definition}

\begin{proposition}\label{prop:gideon-number}
\begin{enumerate}
\item $\gamma(X)=0$ if and only if $X$ is initial, and $\gamma(X)=1$ if and
  only if $X$ is connected.
\item If $\cC$ has small coproducts and the $X_i$ are Gideon-decomposable,
  then $\gamma(\coprod_iX_i)=\sum_i\gamma(X_i)$.
\item $\gamma(X)=|X|$ for a set $X$; $\gamma(X)$ is the number of connected
  components of a space whose components are open; and, in $\Fel(\cC)$, the
  Gideon number of a cat is its number of legs.
\end{enumerate}
\end{proposition}

\begin{proof}
(i) A legless Gideon Decomposition exhibits $X$ as an empty coproduct; a unipedal
one exhibits $X$ as isomorphic to a connected object; conversely an initial
object has the legless Gideon Decomposition and a connected object $X$ has
$(\iota X,\id)$. (ii) Coproducts of Gideon Decompositions are Gideon
Decompositions, by \cref{prop:chloe-basic}(iii) and the associativity of
coproducts. (iii) follows from \cref{ex:gideon}.
\end{proof}

\section{Recognition of categories of cats}\label{sec:recognition}

Let $\cE$ be a locally small category with small coproducts. Then $\cE$ has an
initial object (the empty coproduct). Recall that $\cE^\conn$ is the full
subcategory of connected objects and $y^\conn\colon\cE\to\PSh(\cE^\conn)$ the
restricted Yoneda functor. As noted after \cref{def:gideon}, the notions of
$\cE$-cat and of Gideon Decomposition in $\cE$ make sense without assuming a
terminal object, but there are no $\cE$-cats unless $\cE$ has one. Consequently,
if some object of $\cE$ admits a Gideon Decomposition, or is the vertex of a
cocone from an $\cE$-cat, then $\cE$ has a terminal object.

\begin{theorem}[Recognition]\label{thm:recognition}
For a locally small category $\cE$ the following are equivalent.
\begin{enumerate}
\item\label{it:rec-fel} $\cE\eqv\Fel(\cD)$ for some category $\cD$ with an
  initial and a terminal object.
\item\label{it:rec-gideon} $\cE$ has small coproducts, $\cE$ admits Gideon
  Decompositions, and $\cE^\conn$ has an initial object.
\item\label{it:rec-margaret} $\cE$ has small coproducts, $\cE^\conn$ has an
  initial object, the restricted Yoneda functor
  $y^\conn\colon\cE\to\PSh(\cE^\conn)$ is conservative, and every object of $\cE$
  is the vertex of a Margaret compatible cocone from an $\cE$-cat satisfying the
  Chloe Condition.
\end{enumerate}
In this case $\cD\eqv\cE^\conn$, and $\cE$ has a terminal object.
\end{theorem}

\begin{proof}
\ref{it:rec-fel}$\Rightarrow$\ref{it:rec-gideon}. Each property in
\ref{it:rec-gideon} is invariant under equivalence, so we may take
$\cE=\Fel(\cD)$. It has small coproducts (\cref{prop:coproducts}) and the
terminal object $\Caudal$ (\cref{prop:terminal-cat}); it admits Gideon
Decompositions (\cref{ex:gideon}(v)); and $\Fel(\cD)^\conn$ is the full
subcategory of unipedal cats (\cref{prop:connected-unipedal}), which is
equivalent to $\cD$ via $\iota$ and so has the initial object $\Cephalic$.

\ref{it:rec-gideon}$\Rightarrow$\ref{it:rec-fel}. Let $\cD=\cE^\conn$. By the
remarks preceding the theorem, $\cE$ has a terminal object, so
\cref{prop:chloe-body} applies, and by \cref{prop:chloe-body}(iv),
$\body\colon\Fam(\cD)\to\cE$ is an equivalence. By hypothesis $\cD$ has an
initial object $0_\cD$. The terminal object of $\cE$
corresponds to a terminal object $(S,Z)$ of $\Fam(\cD)$. Since $\cD$ has an
initial object, $\Fam(\cD)((1,0_\cD),(S,Z))\iso S$ as in
\cref{lem:legs-representable}, and this set is a singleton; so $S=1$, and
$Z$ is terminal in $\cD$ because $\cD(W,Z)\iso\Fam(\cD)((1,W),(1,Z))$ is a
singleton for every $W$. Thus $\cD$ has an initial and a terminal object, and
$\cE\eqv\Fam(\cD)\eqv\Fel(\cD)$ by \cref{thm:fundamental}.

\ref{it:rec-gideon}$\Rightarrow$\ref{it:rec-margaret}. By \cref{thm:cmg}\ref{it:cmg-forward},
a Gideon Decomposition of $X$ is a Margaret compatible cocone from a Chloe
$\cE$-cat. For conservativity: by \cref{prop:chloe-body}, every object is
isomorphic to $\body K$ for a Chloe cat $K$, and $y^\conn(\body K)\iso\Yon^\conn K$
naturally in $K$, by connectedness of the objects of $\cE^\conn$. By
\cref{prop:restricted-yoneda,prop:chloe-body}(i), $y^\conn$ is therefore fully
faithful, and in particular conservative.

\ref{it:rec-margaret}$\Rightarrow$\ref{it:rec-gideon}. As before, $\cE$ has a
terminal object. Let $X$ be an object and
$\beta\colon K\to\iota X$ a Margaret compatible cocone with $K$ Chloe. Let
$\beta^\flat\colon\body K\to X$ be its transpose. For connected $Y$,
\[
\coprod_\alpha\cE(Y,L_\alpha)\xrightarrow{\ \kappa\ }\cE(Y,\body K)\xrightarrow{\ \beta^\flat\circ-\ }\cE(Y,X)
\]
composes to $m^\beta_Y$, which is bijective, and $\kappa$ is bijective by
connectedness of $Y$; hence $y^\conn(\beta^\flat)$ is an isomorphism. Since
$y^\conn$ is conservative, $\beta^\flat$ is an isomorphism, and $(K,\beta)$ is a
Gideon Decomposition.

The final assertion follows from \cref{thm:reconstruction} or directly from the
proof of \ref{it:rec-gideon}$\Rightarrow$\ref{it:rec-fel}.
\end{proof}

\begin{remark}\label{rem:recognition-honest}
The equivalence of \ref{it:rec-fel} and \ref{it:rec-gideon}, with $\Fam$ in place
of $\Fel$ and without the conditions on initial and terminal objects, is the
well-known and elementary characterization of free coproduct completions: a
category with small coproducts is a free coproduct completion if and only if
every object is a coproduct of connected objects. We have given a complete
proof and have not attempted to trace the statement to its origin. The
additional conditions record the existence of a head (an initial connected
object) and of a tail (a terminal object, which is built into the notion of a
Gideon Decomposition).
\end{remark}

\begin{counterexample}[Independence of the conditions]\label{cex:recognition}
Each of the three conditions in \cref{thm:recognition}\ref{it:rec-gideon} is
independent of the other two. In each example below, all conditions but the
named one hold.
\begin{enumerate}
\item \emph{Small coproducts.} $\cE=\Kit(\cC)$. The leg functor
  $\legs\colon\Kit(\cC)\to\mathbf{FinSet}$ is left adjoint to $\Delta_1$, so it
  preserves every coproduct that exists in $\Kit(\cC)$. Hence a family of kittens
  has a coproduct in $\Kit(\cC)$ only if all but finitely many of its members are
  legless, and in that case the coproduct of \cref{prop:coproducts} is a kitten
  and is a coproduct in $\Kit(\cC)$. In particular $\Kit(\cC)$ has finite
  coproducts but lacks the coproduct of countably many copies of $\Caudal$, which
  would be sent to a coproduct of countably many copies of $1$ in
  $\mathbf{FinSet}$. With coproducts so described, the proof of
  \cref{prop:connected-unipedal} shows that the connected objects of $\Kit(\cC)$
  are the unipedal kittens, which have the initial object $\Cephalic$; and every
  kitten is a finite coproduct of unipedal kittens, so $\Kit(\cC)$ admits Gideon
  Decompositions.
\item \emph{Initial connected object.} $\cE=\Graph$. It has small coproducts, the
  terminal object consisting of one vertex with a loop, and Gideon
  Decompositions. A connected graph $G_0$ admitting a morphism to the one-vertex
  graph without loop has no edges, hence is a single vertex, and then it admits
  two morphisms to the connected graph with two vertices and one edge. So
  $\Graph^\conn$ has no initial object.
\item \emph{Gideon Decompositions.} $\cE=\Fel(\cD)\times[1]$, where $[1]=\{0<1\}$
  and $\cD$ has initial and terminal objects. Coproducts are computed
  componentwise (in $[1]$ they are suprema), and $(\Caudal,1)$ is terminal. One
  checks from the definitions that the connected objects are the objects
  isomorphic to $(\iota X,0)$: for these,
  $\cE((\iota X,0),\coprod_i(K_i,e_i))\iso\Fel(\cD)(\iota X,\coprod_iK_i)\iso\coprod_i\cE((\iota X,0),(K_i,e_i))$;
  if $K$ is not legless, the identity of $(K,1)\iso(K,0)\sqcup(\Legless,1)$
  factors through neither coprojection, because $\cE((K,1),(K,0))=\varnothing$
  (as $1\not\le0$) and $\cE((K,1),(\Legless,1))=\varnothing$; the object
  $(\Legless,1)$ fails connectedness for the family $((\Legless,1),(\Legless,1))$,
  whose coproduct is $(\Legless,1)$ and which yields two elements on the left of
  \eqref{eq:connected} but one on the right;
  and since $\cE((K,0),(M,e))=\Fel(\cD)(K,M)$ for every $(M,e)$, an object
  $(K,0)$ is connected in $\cE$ if and only if $K$ is connected in $\Fel(\cD)$,
  that is, unipedal (\cref{prop:connected-unipedal}). Thus $\cE^\conn\eqv\cD$ has the initial object
  $(\Cephalic,0)$. But $(K,1)$ is not a coproduct of connected objects, since such
  coproducts have second coordinate $0$.
\end{enumerate}
The last example also shows that conservativity in
\cref{thm:recognition}\ref{it:rec-margaret} cannot be omitted: the cocone from
the Chloe $\cE$-cat $\langle(\iota L_\alpha,0)\rangle_\alpha$ to $(K,1)$, with
components $(j_\alpha,0\le1)$, is Margaret compatible, but it is not a Gideon
Decomposition, and the morphism $(K,0)\to(K,1)$ is inverted by $y^\conn$ without
being an isomorphism.
\end{counterexample}

\begin{remark}\label{rem:recognition-tail}
The category $\Fam(\mathbb N,\le)$ has small coproducts, every object is a
coproduct of connected objects (one-element families), and its connected objects
form a category equivalent to $(\mathbb N,\le)$, which has an initial object. It is
nevertheless not equivalent to any $\Fel(\cD)$, since it has no terminal object.
It does not contradict \cref{thm:recognition}, because, having no terminal object,
it has no cats and hence no Gideon Decompositions. This is one point at which the
requirement that a cat have a tail affects the formulation of a theorem.
\end{remark}

\section{CMG and the tensor product}\label{sec:cmg-tensor}

\begin{proposition}\label{prop:gideon-tensor}
Let $\cC$ be monoidal, with $\otimes$ preserving those small coproducts that exist
in each variable, and suppose that the tensor product of two connected objects
is connected. If $(K,\beta)$ and $(K',\beta')$ are Gideon Decompositions of $X$
and $X'$, then $K\otimes K'$, with the cocone
$(\beta_\alpha\otimes\beta'_{\alpha'})_{(\alpha,\alpha')}$, is a Gideon Decomposition of
$X\otimes X'$. Consequently $\cC^\Gid$ is closed under $\otimes$, and
$\gamma(X\otimes X')=\gamma(X)\gamma(X')$.
\end{proposition}

\begin{proof}
$K\otimes K'$ is Chloe by hypothesis. Preservation of coproducts in each variable
gives $X\otimes X'\iso(\coprod_\alpha L_\alpha)\otimes X'\iso\coprod_\alpha(L_\alpha\otimes\coprod_{\alpha'}L'_{\alpha'})\iso\coprod_{\alpha,\alpha'}L_\alpha\otimes L'_{\alpha'}$,
with coprojections $\beta_\alpha\otimes\beta'_{\alpha'}$.
\end{proof}

\begin{example}\label{ex:gideon-tensor}
The hypotheses of \cref{prop:gideon-tensor} hold for $(\Set,\times)$ and for
$(\Top,\times)$, products of connected spaces being connected and products
distributing over disjoint unions. They fail for $(\Graph,\times)$: the product
of the connected graph $K_2$ (two vertices joined by an edge, no loops) with
itself has vertices $(i,j)$ with $(i,j)\sim(i',j')$ if and only if $i\neq i'$
and $j\neq j'$, and so consists of the two disjoint edges
$\{(0,0),(1,1)\}$ and $\{(0,1),(1,0)\}$. Thus $\gamma(K_2\times K_2)=2\neq1=\gamma(K_2)^2$,
and the conclusion of \cref{prop:gideon-tensor} fails, although $\times$ distributes
over coproducts in $\Graph$.
\end{example}
